
\def\PUDcodigo{ECA206}

\def\PUDcodacad{04507.55}

\def\PUDdisciplina{Robótica II}

\def\PUDsemestre{Opt}

\def\PUDhoras{80}

%NOVOS ---------------------------------------------------
\def\PUDhorasTeoria{40}
\def\PUDhorasPratica{40}

\def\PUDinterECA{ 
\hyperref[PUD:ECA209]{Linguagem de Programação (04507.14)}

\hyperref[PUD:ECA225]{Dispositivos Periféricos (04507.42)} 

\hyperref[PUD:ECA226]{Robótica I (04507.44)}

\hyperref[PUD:EME132]{Inteligência Computacional Aplicada (04507.61)}
}

\def\PUDatuaisECA{- Aplicações com tecnologias atuais.}

\def\PUDrecursos{Projetor multimídia; Quadro branco e pincel.}

%----------------------------------------------------------

\def\PUDcreditos{4}

\def\PUDprerequisito{ \hyperref[PUD:ECA209]{ECA209}  \hyperref[PUD:ECA225]{ECA225} }

\def\PUDpreacad{ 
\hyperref[PUD:ECA226]{Robótica I (04507.44)} 
}

\def\PUDobjetivos{Compreender, projetar e desenvolver sistemas robóticos móveis. Além de verificar o contexto de aplicações regionais e interdiciplinariedade tanto com as disciplinas de pré-requisito quanto com as disicplinas do semestre corrente.
%Compreender as técnicas básicas de modelagem e programação de robôs industriais.  Ler, interpretar e escrever programas na linguagem de robôs industriais.
}

\def\PUDementa{Introdução a Robótica Móvel; locomoção de robôs; Cinemática de robôs móveis; percepção; Visão de máquina aplicada à  Robótica Móvel; localização de robôs  móveis; planejamento e navegação; exemplos de robôs autônomos; aplicações.
%Conceituação de robótica.  Graus de liberdade.  Atuadores e sensores em robótica industrial.  Modelagem e simulação de robôs industriais.  Programação de robôs industriais.  Integração de robôs com periféricos no ambiente industrial.
}

\def\PUDconteudo{ & Introdução à Robótica móvel. Conceitos básicos e aplicações.
%Conceito e definições
\\ 
 & Locomoção. Introdução; Robótica móvel com pernas e com rodas.
 %Tipos de robôs e controladores
\\ 
 & Cinemática em Robótica Móvel. Introdução; restrições e modelos cinemáticos; manobrabilidade; espaço de trabalho e controle de movimento.
 %Tipos de atuadores e sensores
\\ 
 & Percepção. Sensores; Visão Computacional aplicada à Robótica; incerteza na representação e extração de atributos.
 %Métodos de ensino de robôs e de integração com sistemas de automação
\\ 
 & Localização. Introdução; desafios da localização: ruído e \textit{aliasing}; localização baseada em navegação e soluções programadas; representação de crença; representação de mapas; localização probabilística baseada em mapas; sistemas de localização alternativos e construção autônoma de mapas. 
 %Simulação e programação off-line
\\ 
 & Planejamento e navegação. Introdução; competências para navegação: planejamento e reação. Arquiteturas de navegação.
 %Modelagem matemática e simulação
\\ 
 & Inteligência Computacional Aplicada à Robótica. Redes Neurais, Lógica Fuzzy, Algoritmos genéticos, classificadores aplicados à Robótica. Aplicações com tecnologias atuais.
 %Integração com periféricos
\\ 
 & Aplicações. }
 
 Unidade 1 –  Unidade 2 – 
Unidade 3 –  4 – 
Unidade 5 – 
Unidade 6 – 
Unidade 7 – 


\def\PUDmetodo{Aulas expositórias que podem ser teóricas e/ou práticas, onde as práticas no laboratório serão marcadas no decorrer da disciplina.}

\def\PUDavalia{A avaliação será feita com aplicação de provas teóricas e/ou práticas, além da possibilidade de inclusão de trabalhos e seminários no decorrer da disciplina.}

\def\PUDbibbasica{  &  \href{www.biblioteca.ifce.edu.br/index.asp?codigo_sophia=62047}{Niku}, Saeed Benjamin.  Introdução à robótica - análise, controle, aplicações.  2º ed.  LTC, 2015.  382p.  ISBN: 9788521622376.
\\ 
 &  \href{www.biblioteca.ifce.edu.br/index.asp?codigo_sophia=23877}{Rosário}, João Maurício.  Princípios de mecatrônica.  São Paulo - SP.  Pearson Prentice Hall, 2005.  356p.  ISBN: 9788576050100. 
%&  Braunl, Thomas.  Embedded robotics : mobile robot design and applications with embedded systems.  3º ed.  Perth, Austrália : Springer, 2008.  541p.  ISBN 9783540705338
\\ 
&  \href{www.biblioteca.ifce.edu.br/index.asp?codigo_sophia=41222}{Craig}, John J.  Robótica.  3º ed.  São Paulo, SP : Pearson, 2013.  379p.  ISBN: 9788581431284. }

\def\PUDbibcomplementar{ &  \href{www.biblioteca.ifce.edu.br/index.asp?codigo_sophia=62087}{Affonso}, Luiz Otávio Amaral.  Equipamentos mecânicos: análise de falhas e solução de problemas.  3º ed.  Rio de Janeiro: Qualitymark, 2014.  387p.  ISBN: 9788541400367. 
\\
 &  \href{www.biblioteca.ifce.edu.br/index.asp?codigo_sophia=24467}{Coppin}, Ben.  Inteligência artificial.  Rio de Janeiro, RJ : LTC, 2012.  636p.  ISBN: 9788521617297.
 %&  Craig, John J.  Robótica.  3º ed.  Biblioteca Virtual Pearson, 2013.  380p.   ISBN 9788581431284
\\ 
 & \href{www.biblioteca.ifce.edu.br/index.asp?codigo_sophia=62912}{Groover}, Mikell P.  Automação Industrial e Sistemas de Manufatura.  3º ed.  São Paulo, SP: Pearson Prentice Hall, 2015.  581p.  ISBN: 9788576058717.
\\
 & \href{www.biblioteca.ifce.edu.br/index.asp?codigo_sophia=46123}{Bräunl}, Thomas.  Embedded robotics: mobile robot design and applications with embedded systems.  3º ed.  Perth, Austrália: Springer, 2008.  541p.  ISBN: 9783540705338.
\\
 & \href{www.biblioteca.ifce.edu.br/index.asp?codigo_sophia=65450}{Russell}, Stuart.  Inteligência artificial.  Rio de Janeiro, RJ: Elsevier, 2013.  988p.  ISBN: 9788535237016. }

\def\PUDautor{Geraldo Luis Bezerra Ramalho}

\def\PUDdata{2013-04-17}

\def\PUDrevisao{3}

\def\PUDrevisor{Antonio Barbosa de Souza Júnior}

\def\PUDdatarevisao{2019-05-15}

\PUDcorpo
